% A paper skeleton: fast functional mutation via ownership.  The sealed
% discipline, the count of one, the example updates, and the boundary of what is
% claimed.
%
% Build, from this directory. Highlighting is minted's, so the build needs shell
% escape and Pygments on PATH:
%   python3 gen-laws.py     # the law frames and counts.tex, after a library change
%   python3 gen-frames.py   # the frames themselves: the quoted ranges, comment-free
%   latexmk -pdf -shell-escape mutation.tex
% Check that every frame still quotes the library, and that counts.tex holds the
% count the library gives:
%   python3 check-listings.py
%
% pdflatex: the toolchain's lualatex cannot load its font loader, so the Unicode
% in the quoted Lean is mapped to font commands by newunicodechar below. A
% character in a quoted range with no mapping fails the build, which is the
% signal to add one.
\documentclass[11pt,a4paper]{article}

\usepackage[T1]{fontenc}
\usepackage[utf8]{inputenc}
\usepackage{lmodern}
\usepackage{amsmath,amssymb,amsthm}
\usepackage{newunicodechar}
% The mappings the quoted ranges and the prose need, one character per line: a
% character in a quoted range with no mapping fails the build, which is the
% signal to add one.
\newunicodechar{×}{\ensuremath{\times}}
\newunicodechar{σ}{\ensuremath{\sigma}}
\newunicodechar{ρ}{\ensuremath{\rho}}
\newunicodechar{Γ}{\ensuremath{\Gamma}}
\newunicodechar{→}{\ensuremath{\to}}
\newunicodechar{←}{\ensuremath{\leftarrow}}
\newunicodechar{∉}{\ensuremath{\notin}}
\newunicodechar{∧}{\ensuremath{\wedge}}
\newunicodechar{∨}{\ensuremath{\vee}}
\newunicodechar{≤}{\ensuremath{\leq}}
\newunicodechar{⟨}{\ensuremath{\langle}}
\newunicodechar{⟩}{\ensuremath{\rangle}}
\newunicodechar{₀}{\textsubscript{0}}
\newunicodechar{₁}{\textsubscript{1}}
\newunicodechar{₂}{\textsubscript{2}}
\usepackage{booktabs}
\usepackage{enumitem}
\usepackage[margin=1in]{geometry}
\usepackage{microtype}
\usepackage{xcolor}
\usepackage{minted}
% gruvbox-light: the most hues of the styles that stay legible on white paper.
\usemintedstyle{gruvbox-light}
\setminted{fontsize=\footnotesize,frame=single,rulecolor=\color{black!30},
           framesep=4pt,breaklines=true,baselinestretch=1.0}
% A quoted frame is written \codeframe[linerange=lo-hi]{../path}: a line range of
% a library file. The frame itself is generated -- the range with its comments
% removed, since a frame is code and not documentation -- into frames/<key>.lean
% by gen-frames.py, and frames/lookup.tex maps the range to that file. A call
% site therefore names the source it quotes and not the frame, and a range with
% no generated frame is a build error rather than a silently missing box. A
% frame is one unbreakable box, so a frame too tall for the space left on a page
% moves to the next one instead of being split across two.
% Generated by gen-frames.py; do not edit by hand.
% Maps a \codeframe call's (path, range) to its generated frame, so a
% call site names the source and not the frame. The frames themselves are
% frames/<key>.lean, comment-free copies of the ranges, checked by
% check-listings.py against the library.

\expandafter\def\csname cf@../Scidonia/FracPerm/Ledger.lean@25-25\endcsname{9750611a7baa}
\expandafter\def\csname cf@../Scidonia/FracPerm/Ledger.lean@28-28\endcsname{5f928bda90e1}
\expandafter\def\csname cf@../Scidonia/FracPerm/Ledger.lean@34-35\endcsname{1913b63860bf}
\expandafter\def\csname cf@../Scidonia/FracPerm/StateVar/Value.lean@84-84\endcsname{dbe771f805da}
\expandafter\def\csname cf@../Scidonia/FracPerm/StateVar/Value.lean@96-97\endcsname{c011a33daef6}
\expandafter\def\csname cf@../Scidonia/FracPerm/StateVar/Value.lean@109-109\endcsname{c6b92c041b00}
\expandafter\def\csname cf@../Scidonia/FracPerm/StateVar/Value.lean@113-113\endcsname{9c6db65f9b82}
\expandafter\def\csname cf@../Scidonia/FracPerm/StateVar/Model.lean@369-370\endcsname{4ea30c66f461}
\expandafter\def\csname cf@../Scidonia/FracPerm/StateVar/Model.lean@423-424\endcsname{583a4eebbeaa}
\expandafter\def\csname cf@../Scidonia/FracPerm/StateVar/Model.lean@522-523\endcsname{bd6437b6df69}
\expandafter\def\csname cf@../Scidonia/FracPerm/StateVar/Model.lean@529-530\endcsname{a0fef9ee5433}
\expandafter\def\csname cf@../Scidonia/FracPerm/StateVar/Model.lean@535-535\endcsname{57d391554422}
\expandafter\def\csname cf@../Scidonia/FracPerm/StateVar/Model.lean@550-551\endcsname{b15f7df298db}
\expandafter\def\csname cf@../Scidonia/FracPerm/StateVar/Model.lean@597-600\endcsname{ef31681ee890}
\expandafter\def\csname cf@../Scidonia/FracPerm/StateVar/Model.lean@607-608\endcsname{e53944a7355a}
\expandafter\def\csname cf@../Scidonia/FracPerm/StateVar/Model.lean@582-587\endcsname{d4cc154e74a3}
\expandafter\def\csname cf@../Scidonia/FracPerm/StateVar/Model.lean@642-643\endcsname{14055dfba24e}
\expandafter\def\csname cf@../Scidonia/FracPerm/StateVar/Model.lean@656-657\endcsname{ba995450fed5}
\expandafter\def\csname cf@../Scidonia/FracPerm/StateVar/Model.lean@667-668\endcsname{611a3eae16f6}
\expandafter\def\csname cf@../Scidonia/FracPerm/StateVar/Heap.lean@102-103\endcsname{5b740228b3d1}
\expandafter\def\csname cf@../Scidonia/FracPerm/StateVar/Heap.lean@118-119\endcsname{9d55df76baec}
\expandafter\def\csname cf@../Scidonia/FracPerm/StateVar/Heap.lean@126-127\endcsname{548b3c12e0e5}
\expandafter\def\csname cf@../Scidonia/FracPerm/StateVar/Heap.lean@135-140\endcsname{f8af60b72496}
\expandafter\def\csname cf@../Scidonia/FracPerm/StateVar/Heap.lean@157-162\endcsname{8f1dc76ac072}
\expandafter\def\csname cf@../Scidonia/FracPerm/StateVar/Heap.lean@168-171\endcsname{ce7f92922125}
\expandafter\def\csname cf@../Scidonia/FracPerm/StateVar/Heap.lean@214-214\endcsname{c875ce15c57e}
\expandafter\def\csname cf@../Scidonia/FracPerm/StateVar/Array.lean@80-81\endcsname{5084cdad01e1}
\expandafter\def\csname cf@../Scidonia/FracPerm/StateVar/Array.lean@91-92\endcsname{486d61359463}
\expandafter\def\csname cf@../Scidonia/FracPerm/StateVar/Array.lean@113-114\endcsname{dda3851818c2}
\expandafter\def\csname cf@../Scidonia/FracPerm/StateVar/Array.lean@120-121\endcsname{8d41205432b3}
\expandafter\def\csname cf@../Scidonia/FracPerm/StateVar/Array.lean@165-166\endcsname{8041063aa6d3}
\expandafter\def\csname cf@../Scidonia/FracPerm/StateVar/Array.lean@194-195\endcsname{6a9ad9c69b16}
\expandafter\def\csname cf@../Scidonia/FracPerm/StateVar/Array.lean@216-220\endcsname{f98f33160e64}
\expandafter\def\csname cf@../Scidonia/FracPerm/StateVar/Universe.lean@86-87\endcsname{55bff5c33c9b}
\expandafter\def\csname cf@../Scidonia/FracPerm/StateVar/Universe.lean@103-103\endcsname{e23eb9c0c689}
\expandafter\def\csname cf@../Scidonia/FracPerm/StateVar/Universe.lean@153-154\endcsname{90e6c7c31d1e}
\expandafter\def\csname cf@../Scidonia/FracPerm/StateVar/Universe.lean@159-160\endcsname{6b6a3bfec516}
\expandafter\def\csname cf@../Scidonia/FracPerm/StateVar/Universe.lean@178-181\endcsname{e882d85a6188}
\expandafter\def\csname cf@../Scidonia/FracPerm/StateVar/Universe.lean@185-187\endcsname{e28b3a42614c}
\expandafter\def\csname cf@../Scidonia/FracPerm/StateVar/Universe.lean@193-199\endcsname{1fb0ee566945}
\expandafter\def\csname cf@../Scidonia/FracPerm/StateVar/Universe.lean@273-277\endcsname{960c08503b82}
\expandafter\def\csname cf@../Scidonia/FracPerm/StateVar/Slot.lean@67-67\endcsname{191b5865ef60}
\expandafter\def\csname cf@../Scidonia/FracPerm/StateVar/Slot.lean@73-73\endcsname{2e5a81e5133f}
\expandafter\def\csname cf@../Scidonia/FracPerm/StateVar/Slot.lean@82-82\endcsname{d66728bab821}
\expandafter\def\csname cf@../Scidonia/FracPerm/StateVar/Slot.lean@131-132\endcsname{777bd7ea6843}
\expandafter\def\csname cf@../Scidonia/FracPerm/StateVar/Slot.lean@152-156\endcsname{6f3e387a51d5}
\expandafter\def\csname cf@../Scidonia/FracPerm/StateVar/Slot.lean@169-174\endcsname{dca6a2e07dbf}
\expandafter\def\csname cf@../Scidonia/FracPerm/StateVar/Slot.lean@192-192\endcsname{8b3813d52365}
\expandafter\def\csname cf@../Scidonia/FracPerm/StateVar/Layout5.lean@190-191\endcsname{531d56988e55}
\expandafter\def\csname cf@../Scidonia/FracPerm/StateVar/Layout5.lean@194-195\endcsname{81d23c37433c}
\expandafter\def\csname cf@../Scidonia/FracPerm/StateVar/Layout5.lean@198-199\endcsname{3d4fd270215b}
\expandafter\def\csname cf@../Scidonia/FracPerm/StateVar/Layout5.lean@210-211\endcsname{3b81713f1639}
\expandafter\def\csname cf@../Scidonia/FracPerm/StateVar/Layout5.lean@221-222\endcsname{9381cefb1a86}
\expandafter\def\csname cf@../Scidonia/FracPerm/StateVar/Layout5.lean@270-271\endcsname{12f178d518f8}
\expandafter\def\csname cf@../Scidonia/FracPerm/StateVar/Layout5.lean@474-475\endcsname{2e37ef6f1c71}
\expandafter\def\csname cf@../Scidonia/FracPerm/StateVar/Layout5.lean@511-512\endcsname{851dd834f567}
\expandafter\def\csname cf@../Scidonia/FracPerm/StateVar/Layout5.lean@1027-1028\endcsname{0b9db9893a57}
\expandafter\def\csname cf@../Scidonia/FracPerm/StateVar/Layout5.lean@1055-1056\endcsname{79dee75ddae1}
\expandafter\def\csname cf@../Scidonia/FracPerm/StateVar/Layout5.lean@1243-1246\endcsname{c68bb8fcb258}
\expandafter\def\csname cf@../Scidonia/FracPerm/StateVar/Layout5.lean@1265-1267\endcsname{75d40ce95c38}
\expandafter\def\csname cf@../Scidonia/FracPerm/StateVar/Layout5.lean@1283-1284\endcsname{030cb2e59e4d}
\expandafter\def\csname cf@../Scidonia/FracPerm/StateVar/Layout5.lean@1361-1364\endcsname{0f90aede78e5}
\expandafter\def\csname cf@../Scidonia/FracPerm/StateVar/Layout5.lean@1379-1382\endcsname{7501cb075a09}
\expandafter\def\csname cf@../Scidonia/FracPerm/StateVar/Layout5.lean@1402-1405\endcsname{7f8964e9f0e6}
\expandafter\def\csname cf@../Scidonia/FracPerm/StateVar/Sealed.lean@125-127\endcsname{a150e108c8bc}
\expandafter\def\csname cf@../Scidonia/FracPerm/StateVar/Sealed.lean@132-135\endcsname{fa3821e2108b}
\expandafter\def\csname cf@../Scidonia/FracPerm/StateVar/Sealed.lean@152-153\endcsname{232c9c093968}
\expandafter\def\csname cf@../Scidonia/FracPerm/StateVar/Sealed.lean@157-158\endcsname{50d3a390ba9d}
\expandafter\def\csname cf@../Scidonia/FracPerm/StateVar/Sealed.lean@191-193\endcsname{865f8eda760b}
\expandafter\def\csname cf@../Scidonia/FracPerm/StateVar/Sealed.lean@199-202\endcsname{0af04bcb3c10}
\expandafter\def\csname cf@../Scidonia/FracPerm/StateVar/Sealed.lean@213-217\endcsname{4b2ed661cd73}
\expandafter\def\csname cf@../Scidonia/FracPerm/StateVar/Sealed.lean@234-237\endcsname{9b83be61cae0}
\expandafter\def\csname cf@../Scidonia/FracPerm/StateVar/Sealed.lean@251-253\endcsname{e1e96f07c3d9}
\expandafter\def\csname cf@../Scidonia/FracPerm/StateVar/RefCount.lean@740-741\endcsname{c2bd52ed1e21}
\expandafter\def\csname cf@../Scidonia/FracPerm/StateVar/RefCount.lean@765-765\endcsname{c1c6e144a166}
\expandafter\def\csname cf@../Scidonia/FracPerm/StateVar/RefCount.lean@771-771\endcsname{d9b5c6bdb52b}
\expandafter\def\csname cf@../Scidonia/FracPerm/StateVar/RefCount.lean@779-779\endcsname{fbe0489cba29}
\expandafter\def\csname cf@../Scidonia/FracPerm/StateVar/RefCount.lean@787-788\endcsname{163c10511c27}
\expandafter\def\csname cf@../Scidonia/FracPerm/StateVar/RefCount.lean@804-805\endcsname{aa9dea32f52b}
\expandafter\def\csname cf@../Scidonia/FracPerm/StateVar/RefCount.lean@819-823\endcsname{3352aac86792}
\expandafter\def\csname cf@../Scidonia/FracPerm/Structured/Model.lean@292-293\endcsname{331b6c9d7cb0}
\expandafter\def\csname cf@../Scidonia/FracPerm/Structured/Model.lean@297-298\endcsname{51d76a5567c4}
\expandafter\def\csname cf@../Scidonia/FracPerm/Structured/Model.lean@203-204\endcsname{b6db30baf018}
\expandafter\def\csname cf@../Scidonia/FracPerm/Structured/Model.lean@365-366\endcsname{8d3f6d606d28}
\expandafter\def\csname cf@../Scidonia/FracPerm/Structured/Model.lean@372-373\endcsname{0df1d9c5942f}
\expandafter\def\csname cf@../Scidonia/FracPerm/Structured/Model.lean@415-415\endcsname{4210bba08da7}
\expandafter\def\csname cf@../Scidonia/FracPerm/Structured/Model.lean@430-431\endcsname{469cb37dc3cf}
\expandafter\def\csname cf@../Scidonia/FracPerm/Structured/Model.lean@549-552\endcsname{5aa96b675f45}
\expandafter\def\csname cf@../Scidonia/FracPerm/Structured/Model.lean@559-560\endcsname{08b6b25ac8b0}
\expandafter\def\csname cf@../Scidonia/FracPerm/Structured/Model.lean@463-468\endcsname{4f499282495f}
\expandafter\def\csname cf@../Scidonia/FracPerm/Structured/Model.lean@565-567\endcsname{f03a2399da1a}
\expandafter\def\csname cf@../Scidonia/FracPerm/Structured/Model.lean@572-573\endcsname{3908df01cf88}
\expandafter\def\csname cf@../Scidonia/FracPerm/Structured/Model.lean@586-587\endcsname{9ff29af667c0}
\expandafter\def\csname cf@../ScidoniaTest/FracPerm.lean@96-96\endcsname{26c4a5c56142}
\expandafter\def\csname cf@../ScidoniaTest/FracPerm.lean@110-110\endcsname{c7072cc01773}
\expandafter\def\csname cf@../ScidoniaTest/FracPerm.lean@124-124\endcsname{667c4844b79c}
\expandafter\def\csname cf@../ScidoniaTest/FracPerm.lean@136-136\endcsname{af579ae4aafd}
\expandafter\def\csname cf@../ScidoniaTest/FracPerm.lean@149-149\endcsname{7cca48d0769a}
\expandafter\def\csname cf@../ScidoniaTest/FracPerm.lean@160-161\endcsname{58d6bb9503f6}
\expandafter\def\csname cf@../ScidoniaTest/FracPerm.lean@313-313\endcsname{84be31499770}
\expandafter\def\csname cf@../ScidoniaTest/StateVar.lean@281-281\endcsname{846203002cf7}
\expandafter\def\csname cf@../ScidoniaTest/StateVar.lean@287-287\endcsname{6247b6d47e02}
\expandafter\def\csname cf@../ScidoniaTest/StateVar.lean@293-293\endcsname{8ad7a3493477}
\expandafter\def\csname cf@../ScidoniaTest/Model5.lean@335-336\endcsname{3fb4624b4fe2}
\expandafter\def\csname cf@../ScidoniaTest/Model5.lean@353-353\endcsname{a10fbeeae676}
\expandafter\def\csname cf@../ScidoniaTest/Count.lean@44-44\endcsname{743326bf0f37}
\expandafter\def\csname cf@../ScidoniaTest/Count.lean@56-57\endcsname{ce08ec48a432}
\expandafter\def\csname cf@../ScidoniaTest/Count.lean@72-73\endcsname{6ef5df07669f}
\expandafter\def\csname cf@../ScidoniaTest/Count.lean@125-125\endcsname{a34eea8628ee}
\expandafter\def\csname cf@../ScidoniaTest/Count.lean@133-135\endcsname{018ed4c58639}
\expandafter\def\csname cf@../Scidonia/FracPerm/StateVar/Description.lean@24-46\endcsname{952632ccc0d9}
\expandafter\def\csname cf@../Scidonia/FracPerm/StateVar/Layout5.lean@91-119\endcsname{07c4092830a2}
\expandafter\def\csname cf@../Scidonia/FracPerm/StateVar/Layout5.lean@122-138\endcsname{4a24d4844560}
\expandafter\def\csname cf@../Scidonia/FracPerm/StateVar/Layout5.lean@147-166\endcsname{4f032786efd0}
\expandafter\def\csname cf@../Scidonia/FracPerm/StateVar/Layout5.lean@169-177\endcsname{0feaed59d5eb}
\expandafter\def\csname cf@../Scidonia/FracPerm/Ledger.lean@16-22\endcsname{4eb8cea25657}
\expandafter\def\csname cf@../Scidonia/FracPerm/Hoare.lean@107-132\endcsname{0303bf6cccb7}
\expandafter\def\csname cf@../Scidonia/FracPerm/Hoare.lean@145-152\endcsname{a45fb2b10d38}
\expandafter\def\csname cf@../Scidonia/FracPerm/Hoare.lean@178-184\endcsname{aef4c8295edc}
\expandafter\def\csname cf@../Scidonia/FracPerm/Hoare.lean@229-266\endcsname{8cc85285784b}
\expandafter\def\csname cf@../Scidonia/FracPerm/Counter.lean@22-49\endcsname{079f0ad3e32d}
\expandafter\def\csname cf@../Scidonia/FracPerm/StateVar/Value.lean@46-67\endcsname{bfb6d10b4c7b}
\expandafter\def\csname cf@../Scidonia/FracPerm/StateVar/Value.lean@68-82\endcsname{1f6a83f938ee}
\expandafter\def\csname cf@../Scidonia/FracPerm/StateVar/Class.lean@59-91\endcsname{3128a5f31a98}
\expandafter\def\csname cf@../Scidonia/FracPerm/StateVar/Heap.lean@44-47\endcsname{63b5d37b4f3c}
\expandafter\def\csname cf@../Scidonia/FracPerm/StateVar/Heap.lean@49-68\endcsname{0b25f74df449}
\expandafter\def\csname cf@../Scidonia/FracPerm/StateVar/Heap.lean@70-77\endcsname{6c7df423d965}
\expandafter\def\csname cf@../Scidonia/FracPerm/StateVar/Heap.lean@79-89\endcsname{1f6a8050c45e}
\expandafter\def\csname cf@../Scidonia/FracPerm/StateVar/Surface.lean@38-59\endcsname{2ff677e38d0a}
\expandafter\def\csname cf@../Scidonia/FracPerm/StateVar/Array.lean@57-76\endcsname{4d421a7cc316}
\expandafter\def\csname cf@../Scidonia/FracPerm/StateVar/Universe.lean@111-117\endcsname{55053fa59189}
\expandafter\def\csname cf@../Scidonia/FracPerm/StateVar/Universe.lean@210-240\endcsname{e1a2ef0a6af5}
\expandafter\def\csname cf@../Scidonia/FracPerm/StateVar/Slot.lean@46-63\endcsname{d3e2d380990e}
\expandafter\def\csname cf@../Scidonia/FracPerm/StateVar/Sealed.lean@52-60\endcsname{713c79572283}
\expandafter\def\csname cf@../Scidonia/FracPerm/StateVar/Sealed.lean@66-74\endcsname{9ec521f46635}
\expandafter\def\csname cf@../Scidonia/FracPerm/StateVar/Sealed.lean@83-90\endcsname{de0a942c6e6d}
\expandafter\def\csname cf@../Scidonia/FracPerm/StateVar/Layout5.lean@203-208\endcsname{85a1cf2fd656}
\expandafter\def\csname cf@../Scidonia/FracPerm/StateVar/Layout5.lean@298-308\endcsname{5af8e2fc8c99}
\expandafter\def\csname cf@../Scidonia/FracPerm/StateVar/Layout5.lean@457-471\endcsname{118d1cbd1fc1}
\expandafter\def\csname cf@../Scidonia/FracPerm/StateVar/Layout5.lean@539-555\endcsname{d010bd5a10fa}
\expandafter\def\csname cf@../Scidonia/FracPerm/StateVar/Layout5.lean@679-684\endcsname{61f9875d02fa}
\expandafter\def\csname cf@../Scidonia/FracPerm/StateVar/Layout5.lean@710-729\endcsname{95f71547697b}
\expandafter\def\csname cf@../Scidonia/FracPerm/StateVar/Layout5.lean@772-784\endcsname{20bf19049e08}
\expandafter\def\csname cf@../Scidonia/FracPerm/StateVar/RefCount.lean@713-728\endcsname{9729f794a150}
\expandafter\def\csname cf@../Scidonia/FracPerm/Structured/Surface.lean@127-149\endcsname{c9cdc256f152}
\expandafter\def\csname cf@../Scidonia/FracPerm/Structured/Surface.lean@166-179\endcsname{f8b42168ef1a}
\expandafter\def\csname cf@../Scidonia/FracPerm/Structured/Surface.lean@263-270\endcsname{dc0b62954321}

% counts.tex holds the numbers the papers in this directory state about the
% library and about their own panels. This paper's panel count is generated under
% its own name and read below under \panelcount, so the prose says "panel" and
% the numeral comes from the file; check-listings.py fails when that file is
% stale. The other papers' counts keep their own names.
% Generated from the library by gen-laws.py; do not edit by hand.
% The numbers the paper's prose states about the library and about the
% frames it quotes, one \newcommand each, so no numeral is typed by hand.
% The paper's preamble reads this file, and check-listings.py recomputes
% it from the library and from the paper's own call sites.

\newcommand{\theoremcount}{694}  % public theorems under Scidonia/
\newcommand{\vendoredtheoremcount}{17}  % of them vendored from Lean's own model
\newcommand{\lawcount}{110}  % law statements the appendix reproduces
\newcommand{\panelcount}{37}  % panels the paper quotes
\newcommand{\appendixpanelcount}{32}  % of them in the appendix
\newcommand{\refcountframecount}{7}  % frames of the reference-count section

\let\panelcount\mutationpanelcount
\def\codeframe[linerange=#1-#2]#3{%
  \edef\cfkey{cf@#3@#1-#2}%
  \ifcsname\cfkey\endcsname\else
    \PackageError{paper}{no generated frame for #3:#1-#2}%
      {Run python3 gen-frames.py in Paper/.}\fi
  \edef\cffile{frames/\csname\cfkey\endcsname.lean}%
  \par\smallskip\noindent
  \begin{minipage}{\linewidth}%
    \inputminted{lean}{\cffile}%
  \end{minipage}\par}
% A path is set with \path so that it may break at its slashes.
\usepackage{url}\urlstyle{same}
% A body frame is labelled with the file and lines it quotes:
\newcommand{\sourceline}[3]{\par\smallskip\noindent\textbf{#1}\hfill
  {\footnotesize\itshape\path{#2}, lines #3}\par\nobreak}
% Two sections of this draft are prose rather than quotations.  The marker says
% what belongs in the space and keeps the space visible in the built paper.
\newcommand{\stub}[1]{%
  \par\medskip\noindent
  \fcolorbox{black!40}{black!4}{\parbox{\dimexpr\linewidth-2\fboxsep-2\fboxrule}{#1}}%
  \par\medskip}
\usepackage[colorlinks=true,linkcolor=blue!45!black,citecolor=blue!45!black,urlcolor=blue!55!black]{hyperref}
\usepackage{cleveref}

\theoremstyle{definition}
\newtheorem{definition}{Definition}[section]

\newcommand{\code}[1]{\texttt{#1}}

\title{Fast functional mutation via ownership\\[2pt]
\large The sealed discipline, the count of one, the examples, and the boundary}
\author{Gavin Mendel-Gleason\\\small Scidonia}
\date{2026}

\begin{document}
\maketitle

\section{The problem}\label{sec:problem}

\stub{Placeholder. The problem is stated here, in prose. This draft marks the
space it takes and does not fill it; the case for the solution starts at
\cref{sec:sealed}.}

\section{The solution}\label{sec:solution}

\stub{Placeholder. The solution is stated here, in prose, and the rest of the
draft is the case for it: the discipline that establishes the condition
(\cref{sec:sealed}), the count it produces in the model (\cref{sec:count}), the
updates that work (\cref{sec:examples}), and what is not claimed
(\cref{sec:boundary}).}

\section{The sealed approach in Lean}\label{sec:sealed}

A write needs exclusive access to the object it writes, and reads may be shared;
the discipline that arranges this lives in the types and the proofs of the
program, which the compiler erases.
Its last step is a class whose instances supply almost nothing: a state type joins
by naming the layout it is and proving that it is that layout's state type. The
library defines the read, the write and the encoding of a read result from those
two fields, so an instance author cannot choose any of the three, and a client of
a sealed type has no operation that could hand back a mutable part of the state.

\begin{definition}[sealed state variable]
An instance of \code{SealedStateVar} supplies a layout and an equality, and
nothing else. Its read is the layout's read, its write is the layout's write, and
its encoding of a read result is the library's encoding.
\end{definition}

%!check: class SealedStateVar (σ : Type) where
\sourceline{The sealed class}{Scidonia/FracPerm/StateVar/Sealed.lean}{55--59}
\codeframe[linerange=55-59]{../Scidonia/FracPerm/StateVar/Sealed.lean}

The state types that join the class are the shapes a layout denotes: a leaf of a
closed universe of observations, a binary product of two layouts --- so
arbitrarily many fields by nesting --- or an array of states.

%!check: abbrev State : Layout → Type
\sourceline{The shape of a state}{Scidonia/FracPerm/StateVar/Layout5.lean}{122--136}
\codeframe[linerange=122-136]{../Scidonia/FracPerm/StateVar/Layout5.lean}

A read is a path whose end is a leaf of the observation universe. The proof field
is what makes a read name a leaf and nothing else: a path that ends at a
\code{.struct} or an \code{.arr} has no leaf description, so a whole-array read
cannot be written down.

%!check: structure ObsPath (d : Layout) where
\sourceline{The read path}{Scidonia/FracPerm/StateVar/Layout5.lean}{169--175}
\codeframe[linerange=169-175]{../Scidonia/FracPerm/StateVar/Layout5.lean}

The write is one function at a path. It decides first whether the path addresses
anything, rebuilds each struct cell on the path from its updated component, and,
at an array element, takes the element out, checks it, updates it and puts it
back: taking it out is what keeps the array from holding the element while the
update below it runs.

%!check: def setLeaf {d : Layout} (p : ObsPath d) (x : ObsEl p.obs) (s : State d) :
\sourceline{The path-indexed write}{Scidonia/FracPerm/StateVar/Layout5.lean}{778--782}
\codeframe[linerange=778-782]{../Scidonia/FracPerm/StateVar/Layout5.lean}

The sealed surface is that write with the state transported into its layout and
back, reporting what the layout's write reported.

%!check: def setLeaf {σ : Type} [SealedStateVar σ]
\sourceline{The library's write of a sealed state}{Scidonia/FracPerm/StateVar/Sealed.lean}{85--89}
\codeframe[linerange=85-89]{../Scidonia/FracPerm/StateVar/Sealed.lean}

Every Lean quotation in this paper is a line range of the library rather than a
transcription: the \panelcount{} panels here are generated from the source and
re-checked against it, so a range that has drifted from the declaration it quotes
fails the check rather than misquoting it silently.

\section{Ownership semantics leads to reference count one}\label{sec:count}

The discipline is a statement about authority: who may write, and who may hold a
read. The count that decides which of the two implementations of a write runs is a
count of references, and the bridge between the two is a theorem of the library's
counted model. A store typed at an environment that counts its subject once, at a
configuration in the correspondence, is owned by the account and its subject's
reference is counted once: the typing's count of a \emph{name} becomes the
account's count of a \emph{reference}.

%!check: theorem bridge {ρ : TEnv} {m : Multi} {Γ : Occs} {y : Var} {r : Ref} {i : Nat} {v : Expr}
\sourceline{The bridge}{Scidonia/RC/Correspondence.lean}{305--307}
\codeframe[linerange=305-307]{../Scidonia/RC/Correspondence.lean}

At that count, the update is destructive: the written slot keeps its location and
holds what was written, and the cell the store replaced is the cell it wrote. It
is the branch the update takes at a count of one; whether the runtime's own
exclusivity test agrees at that count is the hypothesis \cref{sec:boundary}
names.

%!check: theorem store_owned_is_destructive {c : CHeap} {hs : List Ref} {ρ : TEnv} {m : Multi}
\sourceline{Count one: the store keeps the location}{Scidonia/RC/Ownership.lean}{559--564}
\codeframe[linerange=559-564]{../Scidonia/RC/Ownership.lean}

At any other count it copies, and the theorem states the copy and not only the
loss: the written slots land at a location the heap did not list, the old cell
keeps its slots, and its count drops by one.

%!check: theorem store_shared_copies {c : CHeap} {hs : List Ref} {ρ : TEnv} {m : Multi} {y : Var}
\sourceline{Any other count: the store copies}{Scidonia/RC/Ownership.lean}{600--606}
\codeframe[linerange=600-606]{../Scidonia/RC/Ownership.lean}

So the chain in the model is: a sealed surface that cannot hand out an alias, a
discipline that reaches the write with the object unshared, a bridge that turns
the account's count of the subject's name into one of the subject's reference, and
the destructive branch at that count. \Cref{sec:boundary} says what that chain
does and does not say about the machine.

\section{It actually works: the examples}\label{sec:examples}

Three updates ship with the library, and each is an instance of the same
discipline rather than an arrangement of its own. The first is the flat array: the
primitive's own write, with no discipline around it; the updates below reduce to
it. An index past the end is a caller error rather than
a no-op --- compiled, \code{Array.set!} panics on it --- and the examples write
only to slots the array has.

%!check: def arrSet (a : Array UInt64) (i : Nat) (v : UInt64) : Array UInt64 :=
\sourceline{The flat array}{ScidoniaExamples/Updates.lean}{76--77}
\codeframe[linerange=76-77]{../ScidoniaExamples/Updates.lean}

The second is a nested array. The type is the layout's state type definitionally,
so the instance's \code{rep} is \code{rfl}, and the cast a sealed write performs
erases; the path is two indices, and the take-out, the checks and the put-back at
both levels are the library's.

%!check: def nestedSet (a : Array (Array UInt64)) (i j : Nat) (v : UInt64) :
\sourceline{The nested array}{ScidoniaExamples/Updates.lean}{98--100}
\codeframe[linerange=98-100]{../ScidoniaExamples/Updates.lean}

The third is a family of shapes rather than one shape: the full binary tree of
structs at depth $k$, whose layout is a recursion and whose state type is a nested
pair. A write at a leaf walks $k$ struct cells and rewrites one word.

%!check: def structTree : Nat → Layout5.Layout
\sourceline{The struct tree}{ScidoniaExamples/Updates.lean}{121--123}
\codeframe[linerange=121-123]{../ScidoniaExamples/Updates.lean}

%!check: def treeSet (k : Nat) (v : UInt64) (t : Layout5.State (structTree k)) :
\sourceline{A write at a leaf}{ScidoniaExamples/Updates.lean}{194--196}
\codeframe[linerange=194-196]{../ScidoniaExamples/Updates.lean}

Compiled, the array's two shapes differ by the whole buffer, and the harness
measures both in one program. The exclusive shape carries the array as a recursive
call's parameter, so the only reference the update sees is the callee's; the
aliased shape holds a live second reference across each update, so each one
replaces the buffer. Before the timed runs, each shape is probed once and the
address the update kept or moved is reported beside the times.

%!check: def runBenchmark (seed : Nat) : IO Unit := do
\sourceline{The benchmark}{ScidoniaInPlace/Main.lean}{1814--1825}
\codeframe[linerange=1814-1825]{../ScidoniaInPlace/Main.lean}

\textbf{Measured.} The two tables below are one run of \code{nix develop -c lake
exe inplace} on the machine used here. The times are means over $100$ updates and
$5$ timed measurements per row for the array, and over $200$ updates and $11$
measurements per row for the tree; each row's untimed probes reported what the
columns claim, the exclusive shape keeping the buffer's address and the aliased
shape moving it. Two runs of the harness on this machine differ in the
per-operation times --- most of all in the tree's extra column, which is a
difference of two larger numbers --- so these numerals are one run's, and they are
read for the shape of the columns rather than for the last digit.

\begin{table}[ht]
\centering
\small
\caption{The array update, measured: nanoseconds per update at four sizes. The
times are a measurement of this program on this toolchain, not a theorem, and the
ratio is the ratio of the two measured means.}
\label{tab:bench-array}
\begin{tabular}{rrrr}
\toprule
$n$ & exclusive (ns/op) & aliased (ns/op) & ratio \\
\midrule
$1\,000$ & 30 & 1\,227 & 40.9 \\
$10\,000$ & 30 & 15\,285 & 509.5 \\
$100\,000$ & 41 & 218\,500 & 5\,329.2 \\
$1\,000\,000$ & 41 & 3\,601\,979 & 87\,853.1 \\
\bottomrule
\end{tabular}
\end{table}

\begin{table}[ht]
\centering
\small
\caption{The tree update, measured: nanoseconds per update at four depths, from
the same run and over the same means as \cref{tab:bench-array}. The walk is in
both columns, so the extra column --- the difference of the two --- is the
allocation of the one cell per level the update rebuilds. Measured, not a
theorem.}
\label{tab:bench-tree}
\begin{tabular}{rrrrr}
\toprule
depth & leaves & exclusive (ns/op) & aliased (ns/op) & extra (ns/op) \\
\midrule
8 & 256 & 589 & 668 & 79 \\
12 & 4\,096 & 1\,194 & 1\,246 & 52 \\
16 & 65\,536 & 2\,141 & 2\,197 & 56 \\
20 & 1\,048\,576 & 2\,743 & 3\,003 & 260 \\
\bottomrule
\end{tabular}
\end{table}

The array's aliased column is the copy: it grows with the buffer, from about a
microsecond per update at a thousand elements to $3.6$ milliseconds at a million,
while the exclusive column is flat across the four sizes, at tens of nanoseconds,
and the last row's ratio is above eighty thousand.

The tree's two columns both grow with the depth, because each update walks the
spine either way, and the extra column isolates the allocation: tens to a few
hundred nanoseconds per update. Neither column is a constant-time operation, and
nothing in these four rows separates the walk from the allocation by a theorem;
that separation is in the model, as the allocation of the spine, and
\cref{sec:boundary} states it as the scoped claim it is.

The three updates above are one write each. A fourth example is an algorithm
rather than an update, a union-find, and it is the canonical place where updating
in place rather than copying is the whole point of the program: a union's
amortized cost is near-constant because the two arrays are updated where they lie,
and a copy on every write destroys that. Its arrays are flat arrays of machine
words, sealed the same way the nested array is --- the layout
\code{.arr (.leaf .uint)} has that type as its state type, so the instance's
\code{rep} is \code{rfl} and the cast a sealed write performs erases:

%!check: instance : SealedStateVar (Array UInt64) := ⟨sealedWords, rfl⟩
\sourceline{The union-find's arrays, sealed}{ScidoniaExamples/UnionFind.lean}{70--74}
\codeframe[linerange=70-74]{../ScidoniaExamples/UnionFind.lean}

Each write is therefore the class's take-out update at a one-step path rather than
a bare \code{Array.set!} --- the library's owned exchange, its two checks and its
put-back, with no \code{Layout5} written by hand:

%!check: def writeWord (a : Array UInt64) (i : Nat) (v : UInt64) : Array UInt64 :=
\sourceline{One slot of a machine-word array}{ScidoniaExamples/UnionFind.lean}{89--90}
\codeframe[linerange=89-90]{../ScidoniaExamples/UnionFind.lean}

\code{find} walks from an element to its root and halves the path on the way, so
the walk itself writes; \code{union} finds both roots and links the lower-ranked
one under the higher, which is one write to \code{parent} and, at equal ranks, a
second to \code{rank}:

%!check: def find : Nat → Array UInt64 → Nat → Array UInt64 × Nat
\sourceline{The walk, halving the path}{ScidoniaExamples/UnionFind.lean}{125--132}
\codeframe[linerange=125-132]{../ScidoniaExamples/UnionFind.lean}

%!check: def union (fuel : Nat) (parent rank : Array UInt64) (x y : Nat) :
\sourceline{The link, by rank}{ScidoniaExamples/UnionFind.lean}{143--153}
\codeframe[linerange=143-153]{../ScidoniaExamples/UnionFind.lean}

Two runs perform the same unions of the same fixed sequence on the same two
arrays, and they differ in one thing: what else refers to those arrays while a
step's writes run. The disciplined run carries the arrays as the recursion's
parameters and holds nothing else; the undisciplined run takes a second reference
to both at every step and reads it after that step's writes, which is what makes
the step's first write meet a buffer another name reaches:

%!check: def unionStepsAliased : Nat → List (Nat × Nat) → Array UInt64 → Array UInt64 →
\sourceline{The undisciplined shape}{ScidoniaInPlace/Main.lean}{1298--1310}
\codeframe[linerange=1298-1310]{../ScidoniaInPlace/Main.lean}

The two compute the same partition --- the semantics does not change, only the cost
--- and that is checked rather than assumed: \code{ScidoniaTest/UnionFind}
evaluates both runs on eight elements and four unions and compares the
representatives \code{find} reports, asserts that every reported representative is
its own parent, and states the sealed class's read-back law and its frame law at
the algorithm's write path.

\textbf{Measured.} The table below is a run of the same harness on the same
machine as the two above, and a run of its own: the array and tree tables were
taken before this algorithm's rows existed, so the three are three runs, each
read for the shape of its columns. The means and the probes are the same kind ---
\code{ufReps} measurements of \code{ufOps} unions per row, each shape probed once
before the timed runs. The sequence's length grows with $n$, which is what makes a
row's \emph{total} readable as a shape as well as its per-union mean.

\begin{table}[ht]
\centering
\small
\caption{The union-find, measured: nanoseconds per union and the mean total of the
row's whole sequence, at four element counts. Measured, not a theorem; one run's
numerals, read for the shape of the columns rather than for the last digit.}
\label{tab:bench-union-find}
\begin{tabular}{rrrrrrr}
\toprule
$n$ & unions & excl.\ total ($\mu$s) & aliased total ($\mu$s) & excl.\ (ns/op) & aliased (ns/op) & ratio \\
\midrule
$256$ & 64 & 17.7 & 83.9 & 277 & 1\,310 & 4.7 \\
$1\,024$ & 256 & 67.8 & 1\,195.2 & 264 & 4\,668 & 17.6 \\
$4\,096$ & 1\,024 & 280.5 & 28\,568.9 & 273 & 27\,899 & 102.1 \\
$16\,384$ & 4\,096 & 1\,013.9 & 298\,075.4 & 247 & 72\,772 & 294.6 \\
\bottomrule
\end{tabular}
\end{table}

The exclusive column is flat across the four sizes --- $247$ to $277$ nanoseconds
per union over a $64$-fold range of $n$ --- and its total grows in step with the
sequence, which is one near-constant cost per union. The aliased column is the
copy: $1\,310$ nanoseconds per union at $n = 256$ against $72\,772$ at
$n = 16\,384$, growing with the buffer, so the totals diverge --- $84$ microseconds
to $298$ milliseconds, a factor of about $3\,550$ where $n$ grew by a factor of
$64$ --- and the ratio of the two per-union columns reaches $295$. That is the
contrast the example exists to show: the same algorithm, the same unions, the same
partition, and one near-constant cost against a linear one, decided by nothing but
what else refers to the arrays.

What the discipline gives this algorithm is the sealed take-out update for its
writes and observation-bounded reads --- the class's \code{read?} returns an
element of the closed observation universe, never the array --- and what it does
not give is a proof that any of those writes went in place. That is measured, here
and in \cref{sec:boundary}. The ownership typing the library develops elsewhere
has first-order expressions as its subject and does not reach this Lean function:
the union-find is written \emph{with} the sealed class and is not checked by that
typing.

Two checked children arm the in-place check on this path, and they stop at
different guards. \code{uf-checked-shared} runs the aliased shape and must end the
process naming the \emph{array} site, which is the copy the step would take --- the
loud failure the other examples' children have, and the one the check exists for.
\code{uf-checked-focus} runs the disciplined shape and stops too, at the
\emph{focus} site, and that is a fact about the check rather than about the
update: the array is exclusively owned there, and the element the take-out
extracted is a machine word --- an element of an \code{Array UInt64} is a value and
not a heap object --- for which the guard's exclusivity test answers \code{false}.
The same run with the check off settles it: \code{uf-unique} asserts that both
buffers keep their addresses, so nothing copies. The array guard is therefore what
this path can be armed against, and the in-place evidence for the exclusive run is
that address assertion.

\section{The honest boundary}\label{sec:boundary}

Three seams separate what is proved here from what the machine does.

\begin{enumerate}[leftmargin=1.4em,itemsep=2pt]
\item \emph{The count is the model's.} \Cref{sec:count} proves that a store at
the model's count of one keeps its location, and that any other count copies. The
proof is about the counted heap the library defines and the account the discipline
produces; it is not a statement about Lean's object graph.
\item \emph{The model's count is not the machine's test.} For a compiled write to
be destructive, the runtime's own exclusivity test must accept the object, and the
step from a model count of one to that test is a hypothesis the library names
\code{RuntimeAgrees}. It is stated rather than assumed silently, and no instance
of it is constructed, here or in the checks that accompany the library.

%!check: def RuntimeAgrees (h : Presentation) (handles : List Ref) (r : Ref) (rc : Int32) : Prop :=
\sourceline{The runtime agreement hypothesis}{Scidonia/FracPerm/StateVar/Footprint.lean}{529--530}
\codeframe[linerange=529-530]{../Scidonia/FracPerm/StateVar/Footprint.lean}

\item \emph{In-place update is measured.} The evidence that a compiled write
updates in place is that an executable observed the address the write kept, in
the cases the harness names and in the probes behind \cref{tab:bench-array} and
\cref{tab:bench-union-find}. That is an observation of named cases on one
toolchain, and no theorem here certifies it. The check that turns such a copy into
a terminated process has a limit of its own, which the union-find's cases report:
it tests the exclusivity of the array \emph{and} of the element a take-out
extracted, and a focused element that is a machine word rather than a heap object
stops that test where no copy would happen, so a checked child on such a path
stops at the focus site and the in-place evidence there is the address assertion.
\end{enumerate}

The complexity claim is scoped the same way. For the tree, an update walks
\code{depth} struct cells and rewrites one word in both columns, so the traversal
is $O(\code{depth})$ however the write behaves; what ownership decides is the
\emph{allocation} of the spine --- $O(1)$ cells rebuilt when the tree is
exclusively owned, $O(\log n)$ when a live reference makes it shared. The result
is not an $O(1)$ deep-tree operation.

One limit of the surface belongs beside these, because it bounds what a reader may
conclude from \cref{sec:sealed}: a type joins the sealed class only when it is
definitionally the layout's state type, since the class's \code{rep} is an
equality and not an isomorphism. A user-declared structure would need the
isomorphism instead, and what a conversion between two representations costs has
neither been implemented nor measured.

\end{document}
